\chapter{Unit 3: Data Structures, Classes, and Inheritance}

\section{Section A: Multiple Choice Practice Questions (60 Curated Problems)}

\noindent\textbf{Q1.} Which brackets are commonly used to create a Python list? \hfill [Tag: Lists | Easy]
\begin{enumerate}[label=(\alph*),topsep=2pt,itemsep=1pt,partopsep=0pt,parsep=0pt]
    \item {\{ \}}
    \item {[ ]}
    \item {( )}
    \item {$<$ $>$}
\end{enumerate}
\vspace{2mm}
\noindent\textbf{Q2.} What does the following code produce? [x * 2 for x in range(3)] \hfill [Tag: List comprehension | Easy]
\begin{enumerate}[label=(\alph*),topsep=2pt,itemsep=1pt,partopsep=0pt,parsep=0pt]
    \item {[0, 2, 4]}
    \item {[2, 4, 6]}
    \item {[1, 2, 3]}
    \item {[0, 1, 2]}
\end{enumerate}
\vspace{2mm}
\noindent\textbf{Q3.} Which statement about a Python tuple is correct? \hfill [Tag: Tuples | Easy]
\begin{enumerate}[label=(\alph*),topsep=2pt,itemsep=1pt,partopsep=0pt,parsep=0pt]
    \item {It is immutable}
    \item {It requires square brackets}
    \item {It cannot contain strings}
    \item {It stores only numbers}
\end{enumerate}
\vspace{2mm}
\noindent\textbf{Q4.} Which data structure stores values using key-value pairs? \hfill [Tag: Dictionaries | Easy]
\begin{enumerate}[label=(\alph*),topsep=2pt,itemsep=1pt,partopsep=0pt,parsep=0pt]
    \item {Dictionary}
    \item {List}
    \item {Tuple}
    \item {Set}
\end{enumerate}
\vspace{2mm}
\noindent\textbf{Q5.} What is a key feature of a Python set? \hfill [Tag: Sets: creation, manipulation, and mutability | Easy]
\begin{enumerate}[label=(\alph*),topsep=2pt,itemsep=1pt,partopsep=0pt,parsep=0pt]
    \item {It stores duplicate values}
    \item {It stores unique values}
    \item {It uses numeric indexes}
    \item {It is always immutable}
\end{enumerate}
\vspace{2mm}
\noindent\textbf{Q6.} What is the main purpose of the map() function? \hfill [Tag: Map | Easy]
\begin{enumerate}[label=(\alph*),topsep=2pt,itemsep=1pt,partopsep=0pt,parsep=0pt]
    \item {To sort a sequence}
    \item {To combine two dictionaries}
    \item {To apply a function to items}
    \item {To remove duplicate items}
\end{enumerate}
\vspace{2mm}
\noindent\textbf{Q7.} What does the filter() function do? \hfill [Tag: Filter | Easy]
\begin{enumerate}[label=(\alph*),topsep=2pt,itemsep=1pt,partopsep=0pt,parsep=0pt]
    \item {Creates a class automatically}
    \item {Reverses every item}
    \item {Selects items meeting a condition}
    \item {Combines items into pairs}
\end{enumerate}
\vspace{2mm}
\noindent\textbf{Q8.} What is the usual purpose of reduce() ? \hfill [Tag: Reduce | Easy]
\begin{enumerate}[label=(\alph*),topsep=2pt,itemsep=1pt,partopsep=0pt,parsep=0pt]
    \item {To combine items into one result}
    \item {To remove all strings}
    \item {To generate dictionary keys}
    \item {To create several classes}
\end{enumerate}
\vspace{2mm}
\noindent\textbf{Q9.} What is a Cartesian product of two sets? \hfill [Tag: Cartesian product | Easy]
\begin{enumerate}[label=(\alph*),topsep=2pt,itemsep=1pt,partopsep=0pt,parsep=0pt]
    \item {Their elements in reverse order}
    \item {Only the largest elements}
    \item {All ordered pairs from both sets}
    \item {Their common elements}
\end{enumerate}
\vspace{2mm}
\noindent\textbf{Q10.} What is the purpose of Python's itertools module? \hfill [Tag: Itertools | Easy]
\begin{enumerate}[label=(\alph*),topsep=2pt,itemsep=1pt,partopsep=0pt,parsep=0pt]
    \item {Drawing graphical windows}
    \item {Connecting to databases}
    \item {Changing file permissions}
    \item {Working with iterators}
\end{enumerate}
\vspace{2mm}
\noindent\textbf{Q11.} What does Python's eval() function generally do? \hfill [Tag: eval | Easy]
\begin{enumerate}[label=(\alph*),topsep=2pt,itemsep=1pt,partopsep=0pt,parsep=0pt]
    \item {Defines a new module}
    \item {Deletes a string variable}
    \item {Prints every list item}
    \item {Executes a string as an expression}
\end{enumerate}
\vspace{2mm}
\noindent\textbf{Q12.} What does Python's exec() function generally execute? \hfill [Tag: exec | Easy]
\begin{enumerate}[label=(\alph*),topsep=2pt,itemsep=1pt,partopsep=0pt,parsep=0pt]
    \item {A list of dictionary keys}
    \item {A string containing Python statements}
    \item {Only imported packages}
    \item {Only integer values}
\end{enumerate}
\vspace{2mm}
\noindent\textbf{Q13.} What does enumerate() add when iterating over a sequence? \hfill [Tag: enumerate | Easy]
\begin{enumerate}[label=(\alph*),topsep=2pt,itemsep=1pt,partopsep=0pt,parsep=0pt]
    \item {Files and folders}
    \item {Keys and methods}
    \item {Indexes and values}
    \item {Classes and objects}
\end{enumerate}
\vspace{2mm}
\noindent\textbf{Q14.} What does the zip() function commonly do? \hfill [Tag: zip | Easy]
\begin{enumerate}[label=(\alph*),topsep=2pt,itemsep=1pt,partopsep=0pt,parsep=0pt]
    \item {Copies a complete class}
    \item {Compresses a Python file}
    \item {Pairs items from iterables}
    \item {Sorts all numbers}
\end{enumerate}
\vspace{2mm}
\noindent\textbf{Q15.} Which module provides common functions for copying objects? \hfill [Tag: copy | Easy]
\begin{enumerate}[label=(\alph*),topsep=2pt,itemsep=1pt,partopsep=0pt,parsep=0pt]
    \item {time}
    \item {math}
    \item {random}
    \item {copy}
\end{enumerate}
\vspace{2mm}
\noindent\textbf{Q16.} What is a class in Python? \hfill [Tag: Introduction to classes | Easy]
\begin{enumerate}[label=(\alph*),topsep=2pt,itemsep=1pt,partopsep=0pt,parsep=0pt]
    \item {A command-line argument}
    \item {A type of loop}
    \item {A blueprint for objects}
    \item {A built-in file format}
\end{enumerate}
\vspace{2mm}
\noindent\textbf{Q17.} Which special method is commonly used as a constructor in a Python class? \hfill [Tag: Constructors | Easy]
\begin{enumerate}[label=(\alph*),topsep=2pt,itemsep=1pt,partopsep=0pt,parsep=0pt]
    \item {\_\_create\_\_()}
    \item {\_\_start\_\_()}
    \item {\_\_newobject\_\_()}
    \item {\_\_init\_\_()}
\end{enumerate}
\vspace{2mm}
\noindent\textbf{Q18.} What is an attribute in an object-oriented Python program? \hfill [Tag: Attributes | Easy]
\begin{enumerate}[label=(\alph*),topsep=2pt,itemsep=1pt,partopsep=0pt,parsep=0pt]
    \item {A special type of comment}
    \item {Data associated with an object}
    \item {A package installation command}
    \item {A repeated conditional statement}
\end{enumerate}
\vspace{2mm}
\noindent\textbf{Q19.} What is method overriding? \hfill [Tag: Method overriding | Easy]
\begin{enumerate}[label=(\alph*),topsep=2pt,itemsep=1pt,partopsep=0pt,parsep=0pt]
    \item {Removing every method in a class}
    \item {Redefining an inherited method}
    \item {Calling a method without parentheses}
    \item {Adding a variable outside Python}
\end{enumerate}
\vspace{2mm}
\noindent\textbf{Q20.} Which keyword is used in a generator function to produce values one at a time? \hfill [Tag: Use of decorators and generators | Easy]
\begin{enumerate}[label=(\alph*),topsep=2pt,itemsep=1pt,partopsep=0pt,parsep=0pt]
    \item {return}
    \item {produce}
    \item {generate}
    \item {yield}
\end{enumerate}
\vspace{2mm}
\noindent\textbf{Q21.} What is the value of nums after this code runs? \hfill [Tag: Lists | Medium]
\begin{enumerate}[label=(\alph*),topsep=2pt,itemsep=1pt,partopsep=0pt,parsep=0pt]
    \item {[1, 7, 3]}
    \item {[1, 7, 8, 3]}
    \item {[1, 2, 7, 8, 3]}
    \item {[7, 8, 1, 2, 3]}
\end{enumerate}
\vspace{2mm}
\noindent\textbf{Q22.} What does the following list comprehension produce? \hfill [Tag: List comprehension | Medium]
\begin{enumerate}[label=(\alph*),topsep=2pt,itemsep=1pt,partopsep=0pt,parsep=0pt]
    \item {[1, 9, 25]}
    \item {[0, 4, 16]}
    \item {[0, 1, 4, 9, 16, 25]}
    \item {[1, 4, 9, 16, 25]}
\end{enumerate}
\vspace{2mm}
\noindent\textbf{Q23.} Which statement correctly describes the result of this code? \hfill [Tag: Tuples | Medium]
\begin{enumerate}[label=(\alph*),topsep=2pt,itemsep=1pt,partopsep=0pt,parsep=0pt]
    \item {The list changes but the tuple structure remains fixed}
    \item {The tuple becomes ([1, 2, 4], 3)}
    \item {The tuple changes into a list automatically}
    \item {The code raises a TypeError}
\end{enumerate}
\vspace{2mm}
\noindent\textbf{Q24.} What is printed by this code? \hfill [Tag: Dictionaries | Medium]
\begin{enumerate}[label=(\alph*),topsep=2pt,itemsep=1pt,partopsep=0pt,parsep=0pt]
    \item {1 \{'a': 1, 'b': 2\}}
    \item {2 \{'b': 2\}}
    \item {0 \{'a': 1\}}
    \item {1 \{'b': 2\}}
\end{enumerate}
\vspace{2mm}
\noindent\textbf{Q25.} What is the value of result after this code? \hfill [Tag: Sets: creation, manipulation, and mutability | Medium]
\begin{enumerate}[label=(\alph*),topsep=2pt,itemsep=1pt,partopsep=0pt,parsep=0pt]
    \item {\{2, 3, 4\}}
    \item {\{1, 4\}}
    \item {\{2, 3\}}
    \item {\{1, 2, 3, 4, 5\}}
\end{enumerate}
\vspace{2mm}
\noindent\textbf{Q26.} What does this expression evaluate to when converted to a list? \hfill [Tag: Map | Medium]
\begin{enumerate}[label=(\alph*),topsep=2pt,itemsep=1pt,partopsep=0pt,parsep=0pt]
    \item {["cat", "python", "AI"]}
    \item {[2, 5, 1]}
    \item {[3, 5, 2]}
    \item {[3, 6, 2]}
\end{enumerate}
\vspace{2mm}
\noindent\textbf{Q27.} What is the result of the following expression? \hfill [Tag: Filter | Medium]
\begin{enumerate}[label=(\alph*),topsep=2pt,itemsep=1pt,partopsep=0pt,parsep=0pt]
    \item {[1, 3, 2]}
    \item {[4, 6]}
    \item {[1, 4, 3, 6, 2]}
    \item {[3, 4, 6]}
\end{enumerate}
\vspace{2mm}
\noindent\textbf{Q28.} What value is produced by this code? \hfill [Tag: Reduce | Medium]
\begin{enumerate}[label=(\alph*),topsep=2pt,itemsep=1pt,partopsep=0pt,parsep=0pt]
    \item {9}
    \item {12}
    \item {10}
    \item {24}
\end{enumerate}
\vspace{2mm}
\noindent\textbf{Q29.} How many pairs are generated by the Cartesian product of these two lists? \hfill [Tag: Cartesian product | Medium]
\begin{enumerate}[label=(\alph*),topsep=2pt,itemsep=1pt,partopsep=0pt,parsep=0pt]
    \item {3}
    \item {6}
    \item {9}
    \item {5}
\end{enumerate}
\vspace{2mm}
\noindent\textbf{Q30.} What does the following expression produce? \hfill [Tag: Itertools | Medium]
\begin{enumerate}[label=(\alph*),topsep=2pt,itemsep=1pt,partopsep=0pt,parsep=0pt]
    \item {[1, 2, [3], 4, 5]}
    \item {[5, 4, 3, 2, 1]}
    \item {[1, 2, 3, 4, 5]}
    \item {[[1, 2], [3], [4, 5]]}
\end{enumerate}
\vspace{2mm}
\noindent\textbf{Q31.} What is printed by this code? \hfill [Tag: eval | Medium]
\begin{enumerate}[label=(\alph*),topsep=2pt,itemsep=1pt,partopsep=0pt,parsep=0pt]
    \item {14}
    \item {A SyntaxError is raised}
    \item {18}
    \item {"3 * (2 + 4)"}
\end{enumerate}
\vspace{2mm}
\noindent\textbf{Q32.} What is the value of total after this code executes? \hfill [Tag: exec | Medium]
\begin{enumerate}[label=(\alph*),topsep=2pt,itemsep=1pt,partopsep=0pt,parsep=0pt]
    \item {9}
    \item {The value is None}
    \item {"4 + 5"}
    \item {An undefined variable error occurs}
\end{enumerate}
\vspace{2mm}
\noindent\textbf{Q33.} What is printed by this code? \hfill [Tag: enumerate | Medium]
\begin{enumerate}[label=(\alph*),topsep=2pt,itemsep=1pt,partopsep=0pt,parsep=0pt]
    \item {1 blue followed by 2 red}
    \item {0 blue followed by 1 red}
    \item {0 red followed by 1 blue}
    \item {1 red followed by 2 blue}
\end{enumerate}
\vspace{2mm}
\noindent\textbf{Q34.} What is the result of this expression? \hfill [Tag: zip | Medium]
\begin{enumerate}[label=(\alph*),topsep=2pt,itemsep=1pt,partopsep=0pt,parsep=0pt]
    \item {[(1, 'a'), (2, 'b')]}
    \item {[(1, 2, 3), ('a', 'b')]}
    \item {[(1, 'a'), (2, 'b'), (3, None)]}
    \item {[(1, 'a'), (2, 'b'), (3, '')]}
\end{enumerate}
\vspace{2mm}
\noindent\textbf{Q35.} What is printed by this code? \hfill [Tag: copy | Medium]
\begin{enumerate}[label=(\alph*),topsep=2pt,itemsep=1pt,partopsep=0pt,parsep=0pt]
    \item {A TypeError is raised}
    \item {[[1, 2, 9], [3, 4]]}
    \item {[[1, 2], [3, 4]]}
    \item {[[9], [3, 4]]}
\end{enumerate}
\vspace{2mm}
\noindent\textbf{Q36.} What is the output of this code? \hfill [Tag: Constructors | Medium]
\begin{enumerate}[label=(\alph*),topsep=2pt,itemsep=1pt,partopsep=0pt,parsep=0pt]
    \item {An AttributeError is raised}
    \item {None}
    \item {0}
    \item {5}
\end{enumerate}
\vspace{2mm}
\noindent\textbf{Q37.} What is printed by this code? \hfill [Tag: Methods | Medium]
\begin{enumerate}[label=(\alph*),topsep=2pt,itemsep=1pt,partopsep=0pt,parsep=0pt]
    \item {125}
    \item {75}
    \item {100}
    \item {25}
\end{enumerate}
\vspace{2mm}
\noindent\textbf{Q38.} What is printed by this code? \hfill [Tag: Attributes | Medium]
\begin{enumerate}[label=(\alph*),topsep=2pt,itemsep=1pt,partopsep=0pt,parsep=0pt]
    \item {North Central}
    \item {North North}
    \item {An AttributeError is raised}
    \item {Central Central}
\end{enumerate}
\vspace{2mm}
\noindent\textbf{Q39.} What is printed by this code? \hfill [Tag: Method overriding | Medium]
\begin{enumerate}[label=(\alph*),topsep=2pt,itemsep=1pt,partopsep=0pt,parsep=0pt]
    \item {Dog}
    \item {A TypeError is raised}
    \item {bark}
    \item {generic}
\end{enumerate}
\vspace{2mm}
\noindent\textbf{Q40.} What is printed by this code? \hfill [Tag: Multiple inheritance | Medium]
\begin{enumerate}[label=(\alph*),topsep=2pt,itemsep=1pt,partopsep=0pt,parsep=0pt]
    \item {C}
    \item {A}
    \item {B}
    \item {A method resolution error occurs}
\end{enumerate}
\vspace{2mm}
\noindent\textbf{Q41.} What is printed by the following code? a = [1, [2, 3]]; b = a[:]; a[1].append(4); b.append(5); print(a, b) \hfill [Tag: Lists | Hard]
\begin{enumerate}[label=(\alph*),topsep=2pt,itemsep=1pt,partopsep=0pt,parsep=0pt]
    \item {[1, [2, 3, 4]] [1, [2, 3]]}
    \item {[1, [2, 3, 4], 5] [1, [2, 3, 4], 5]}
    \item {[1, [2, 3, 4]] [1, [2, 3, 4], 5]}
    \item {[1, [2, 3], 5] [1, [2, 3, 4]]}
\end{enumerate}
\vspace{2mm}
\noindent\textbf{Q42.} What is the value of result after execution? data = [1, 2, 3]; result = [data.pop(0) for \_ in range(len(data))] \hfill [Tag: List comprehension | Hard]
\begin{enumerate}[label=(\alph*),topsep=2pt,itemsep=1pt,partopsep=0pt,parsep=0pt]
    \item {An IndexError is raised}
    \item {[1, 2, 3]}
    \item {[1, 2]}
    \item {[3, 2, 1]}
\end{enumerate}
\vspace{2mm}
\noindent\textbf{Q43.} Which statement correctly explains the result of hash((1, [2, 3])) ? \hfill [Tag: Tuples | Hard]
\begin{enumerate}[label=(\alph*),topsep=2pt,itemsep=1pt,partopsep=0pt,parsep=0pt]
    \item {It raises TypeError because the tuple contains an unhashable list}
    \item {It returns different hashes each time because lists are mutable}
    \item {It returns a stable hash because tuples are immutable}
    \item {It hashes the list by converting it to a tuple automatically}
\end{enumerate}
\vspace{2mm}
\noindent\textbf{Q44.} What is the value of merged ? left = \{"a": 1, "b": 2\}; right = \{"b": 9, "c": 3\}; merged = left | right \hfill [Tag: Dictionaries | Hard]
\begin{enumerate}[label=(\alph*),topsep=2pt,itemsep=1pt,partopsep=0pt,parsep=0pt]
    \item {A KeyError is raised for the duplicate key}
    \item {\{"a": 1, "b": 2, "b": 9, "c": 3\}}
    \item {\{"a": 1, "b": 2, "c": 3\}}
    \item {\{"a": 1, "b": 9, "c": 3\}}
\end{enumerate}
\vspace{2mm}
\noindent\textbf{Q45.} What is the most appropriate behavior of this code? s = \{1, 2, 3\}; for x in s: s.remove(x) \hfill [Tag: Sets: creation, manipulation, and mutability | Hard]
\begin{enumerate}[label=(\alph*),topsep=2pt,itemsep=1pt,partopsep=0pt,parsep=0pt]
    \item {A RuntimeError is raised because the set size changes during iteration}
    \item {Only the first element is removed, then iteration stops normally}
    \item {A new set is created automatically for safe iteration}
    \item {All elements are removed successfully}
\end{enumerate}
\vspace{2mm}
\noindent\textbf{Q46.} What does the following expression evaluate to? from functools import reduce; reduce(lambda a, b: a * b, filter(lambda x: x \% 2, map(lambda x: x + 1, [1, 2, 3, 4])), 1) \hfill [Tag: Map, Filter, Reduce | Hard]
\begin{enumerate}[label=(\alph*),topsep=2pt,itemsep=1pt,partopsep=0pt,parsep=0pt]
    \item {45}
    \item {120}
    \item {15}
    \item {24}
\end{enumerate}
\vspace{2mm}
\noindent\textbf{Q47.} What is the value of list(itertools.product([1, 2], repeat=2)) ? \hfill [Tag: Cartesian product | Hard]
\begin{enumerate}[label=(\alph*),topsep=2pt,itemsep=1pt,partopsep=0pt,parsep=0pt]
    \item {[(1, 2), (2, 1)]}
    \item {[(1, 1), (1, 2), (2, 1), (2, 2)]}
    \item {[(1, 1), (2, 2)]}
    \item {[(1, 1, 2, 2)]}
\end{enumerate}
\vspace{2mm}
\noindent\textbf{Q48.} What is printed by this code? from itertools import groupby; values = [1, 1, 2, 2, 1]; print([(k, list(g)) for k, g in groupby(values)]) \hfill [Tag: Itertools | Hard]
\begin{enumerate}[label=(\alph*),topsep=2pt,itemsep=1pt,partopsep=0pt,parsep=0pt]
    \item {[(1, [1, 1]), (1, [2, 2]), (1, [1])]}
    \item {[(1, [1, 1]), (2, [2, 2]), (1, [1])]}
    \item {[(1, [1]), (1, [1]), (2, [2]), (2, [2]), (1, [1])]}
    \item {[(1, [1, 1, 1]), (2, [2, 2])]}
\end{enumerate}
\vspace{2mm}
\noindent\textbf{Q49.} What happens when this expression is evaluated? eval("\_\_import\_\_('os').getcwd()", \{"\_\_builtins\_\_": \{\}\}, \{\}) \hfill [Tag: eval | Hard]
\begin{enumerate}[label=(\alph*),topsep=2pt,itemsep=1pt,partopsep=0pt,parsep=0pt]
    \item {It raises NameError because \_\_import\_\_ is unavailable}
    \item {It raises SyntaxError because imports are forbidden in eval}
    \item {It silently returns None because built-ins are disabled}
    \item {It returns the current working directory}
\end{enumerate}
\vspace{2mm}
\noindent\textbf{Q50.} What happens when f() is called? g = \{\}; l = \{\}; exec("x = 10 def f(): return x", g, l); f = l["f"]; print(f()) \hfill [Tag: exec | Hard]
\begin{enumerate}[label=(\alph*),topsep=2pt,itemsep=1pt,partopsep=0pt,parsep=0pt]
    \item {It prints 10 because x is in the local execution namespace}
    \item {It raises NameError because the function searches g for global x}
    \item {It prints None because exec does not retain function definitions}
    \item {It raises KeyError because x is not in g}
\end{enumerate}
\vspace{2mm}
\noindent\textbf{Q51.} What is the value of list(enumerate("ab", start=3)) ? \hfill [Tag: enumerate | Hard]
\begin{enumerate}[label=(\alph*),topsep=2pt,itemsep=1pt,partopsep=0pt,parsep=0pt]
    \item {[(0, 'a'), (1, 'b')]}
    \item {[(3, 'b'), (4, 'a')]}
    \item {[(3, 'a'), (4, 'b')]}
    \item {[(1, 'a'), (2, 'b')]}
\end{enumerate}
\vspace{2mm}
\noindent\textbf{Q52.} What is the value of list(zip([1, 2, 3], "xy")) ? \hfill [Tag: zip | Hard]
\begin{enumerate}[label=(\alph*),topsep=2pt,itemsep=1pt,partopsep=0pt,parsep=0pt]
    \item {[(1, 'x'), (2, 'y'), (3, None)]}
    \item {[(1, 'x'), (2, 'y')]}
    \item {[(1, 2, 3), ('x', 'y')]}
    \item {A ValueError is raised because the lengths differ}
\end{enumerate}
\vspace{2mm}
\noindent\textbf{Q53.} What is printed? import copy; a = [[1], [2]]; b = copy.copy(a); c = copy.deepcopy(a); a[0].append(9); print(b, c) \hfill [Tag: copy | Hard]
\begin{enumerate}[label=(\alph*),topsep=2pt,itemsep=1pt,partopsep=0pt,parsep=0pt]
    \item {[[1, 9], [2]] [[1], [2]]}
    \item {[[1, 9], [2]] [[1, 9], [2]]}
    \item {[[1], [2]] [[1], [2]]}
    \item {[[1], [2]] [[1, 9], [2]]}
\end{enumerate}
\vspace{2mm}
\noindent\textbf{Q54.} What is printed? class A: value = 1 a = A(); b = A(); a.value = 2; print(a.value, b.value, A.value) \hfill [Tag: Attributes | Hard]
\begin{enumerate}[label=(\alph*),topsep=2pt,itemsep=1pt,partopsep=0pt,parsep=0pt]
    \item {2 2 2}
    \item {1 2 1}
    \item {1 1 1}
    \item {2 1 1}
\end{enumerate}
\vspace{2mm}
\noindent\textbf{Q55.} What is the result of this code? class Base: def \_\_init\_\_(self): self.x = 1 class Child(Base): def \_\_init\_\_(self): self.y = 2 c = Child(); print(hasattr(c, "x"), c.y) \hfill [Tag: Constructors | Hard]
\begin{enumerate}[label=(\alph*),topsep=2pt,itemsep=1pt,partopsep=0pt,parsep=0pt]
    \item {False 2}
    \item {True 1}
    \item {True 2}
    \item {An AttributeError is raised while constructing c}
\end{enumerate}
\vspace{2mm}
\noindent\textbf{Q56.} What is printed? class A: def show(self): return "A" class B(A): def show(self): return super().show() + "B" print(B().show()) \hfill [Tag: Method overriding | Hard]
\begin{enumerate}[label=(\alph*),topsep=2pt,itemsep=1pt,partopsep=0pt,parsep=0pt]
    \item {BA}
    \item {B}
    \item {A}
    \item {AB}
\end{enumerate}
\vspace{2mm}
\noindent\textbf{Q57.} Given class Animal: pass and class Dog(Animal): pass , which statement is correct? \hfill [Tag: Inheritance concepts | Hard]
\begin{enumerate}[label=(\alph*),topsep=2pt,itemsep=1pt,partopsep=0pt,parsep=0pt]
    \item {isinstance(Dog, Animal) is true}
    \item {issubclass(Dog, Animal) is true}
    \item {isinstance(Animal(), Dog) is true}
    \item {issubclass(Animal, Dog) is true}
\end{enumerate}
\vspace{2mm}
\noindent\textbf{Q58.} What is printed? class A: def f(self): return "A" class B(A): def f(self): return "B" class C(A): def f(self): return "C" class D(B, C): pass print(D().f()) \hfill [Tag: Multiple inheritance | Hard]
\begin{enumerate}[label=(\alph*),topsep=2pt,itemsep=1pt,partopsep=0pt,parsep=0pt]
    \item {A method-resolution conflict is raised}
    \item {A}
    \item {C}
    \item {B}
\end{enumerate}
\vspace{2mm}
\noindent\textbf{Q59.} Which design best preserves encapsulation for a bank account whose balance must never be directly set to a negative value? \hfill [Tag: Practical OOP examples | Hard]
\begin{enumerate}[label=(\alph*),topsep=2pt,itemsep=1pt,partopsep=0pt,parsep=0pt]
    \item {Use a class variable shared by every account}
    \item {Store \_balance and validate changes through methods}
    \item {Expose balance publicly and trust callers}
    \item {Store the balance in a global variable for central control}
\end{enumerate}
\vspace{2mm}
\noindent\textbf{Q60.} What is printed? def deco(f): def wrapper(): return f() * 2 return wrapper @deco def value(): return 3 print(value()) \hfill [Tag: Use of decorators and generators | Hard]
\begin{enumerate}[label=(\alph*),topsep=2pt,itemsep=1pt,partopsep=0pt,parsep=0pt]
    \item {3}
    \item {A TypeError is raised because decorators cannot alter return values}
    \item {5}
    \item {6}
\end{enumerate}
\vspace{2mm}

\section{Section B: Comprehensive Subjective Problem with Detailed Analytical Rubric}

\begin{subjectiveq}{3}{1}
Implement an Object-Oriented Least-Recently-Used (LRU) Cache class \texttt{LRUCache(capacity)} supporting $\mathcal{O}(1)$ average-time \texttt{get(key)} and \texttt{put(key, value)} operations. Employ a custom doubly linked list and a Python dictionary. Demonstrate inheritance by creating a specialized \texttt{ExpiringLRUCache} subclass that supports time-to-live (TTL) expiration.
\end{subjectiveq}

\begin{solbox}
\begin{lstlisting}
import time

class Node:
    def __init__(self, key, val):
        self.key, self.val = key, val
        self.prev = self.next = None

class LRUCache:
    def __init__(self, capacity: int):
        self.capacity = capacity
        self.cache = {}
        self.head, self.tail = Node(0, 0), Node(0, 0)
        self.head.next = self.tail
        self.tail.prev = self.head

    def _remove(self, node):
        node.prev.next = node.next
        node.next.prev = node.prev

    def _add(self, node):
        node.next = self.head.next
        node.prev = self.head
        self.head.next.prev = node
        self.head.next = node

    def get(self, key):
        if key in self.cache:
            node = self.cache[key]
            self._remove(node)
            self._add(node)
            return node.val
        return -1

    def put(self, key, value):
        if key in self.cache:
            self._remove(self.cache[key])
        node = Node(key, value)
        self._add(node)
        self.cache[key] = node
        if len(self.cache) > self.capacity:
            lru = self.tail.prev
            self._remove(lru)
            del self.cache[lru.key]

class ExpiringLRUCache(LRUCache):
    def __init__(self, capacity: int, default_ttl: float):
        super().__init__(capacity)
        self.ttl = default_ttl
        self.timestamps = {}

    def get(self, key):
        if key in self.timestamps and (time.time() - self.timestamps[key]) > self.ttl:
            # Expired: remove and return -1
            self._remove(self.cache[key])
            del self.cache[key]
            del self.timestamps[key]
            return -1
        return super().get(key)
\end{lstlisting}
\textbf{Complexity \& Data Structure Analysis:}\\
$\bullet$ Average operations achieve strict $\mathcal{O}(1)$ time through dictionary key hash lookup and constant-time pointer reassignment in the doubly linked list.
\end{solbox}

\begin{examinertip}
$\bullet$ \textbf{Rubric Breakdown (10 Marks Total):}\\
1. Algorithmic Formulation \& Theoretical Correctness: \textbf{3 Marks}\\
2. Production Python Syntax \& Idiomatic Standards (PEP 8): \textbf{3 Marks}\\
3. Boundary Case \& Exception Robustness: \textbf{2 Marks}\\
4. Asymptotic Time \& Auxiliary Space Complexity Derivation: \textbf{2 Marks}
\end{examinertip}
